\section{My role}

\textbf{1a. What I was given vs what I ended up doing.} \rtag{M16.a} On paper I was the ``CV and autonomy lead.'' That label survived the first kickoff meeting in wk12 and then started drifting almost immediately. By wk16 I owned the state machine (\texttt{main.py}, 1411 lines at peak), the simulation framework (\texttt{simple\_simulator.py}, 2508 lines), the vision pipeline (\texttt{vision.py} dual-backend), the geofence layer, the browser ground station (\texttt{pi\_flight.py}, 1097 lines), the progressive test ladder (41 scripts), the Pi integration, and the project's ground-truth document (\texttt{CLAUDE.md}, 800+ lines). By wk20 I had also taken on teammate-facing scaffolding: the Robin integration script \cite{robinhandoff}, Demetro's Final Design Review (FDR) speaking scripts, the colleague onboarding doc, and by wk22 the team's formal complaint letter to the course organiser \cite{complaintletter}.

The explicit/implicit split was stark. Explicit: CV and autonomy. Implicit and drifted-into: simulation, safety layer, report writing, teammate-enablement, external communication. Our academic supervisor gave feedback at the milestone reviews (PDR, IDR, FDR) but had no day-to-day visibility into how work was being distributed; neither he nor our project manager saw the slow concentration of software ownership happening until it was already a fact. Nobody allocated these to me; I picked them up because the alternative was for them not to exist. My Belbin \cite{belbin2010} profile through the project was \textit{Plant} early (ideas, architecture), \textit{Implementer} in the middle (writing the code that made the ideas real), and \textit{Completer--Finisher} at the end (127 unit tests, modularity audit 3.8/5.0, capability audit Sim~7.7/10 vs Real~2.1/10, R01--R12 gap-closure on 11 April). I was not a \textit{Coordinator} -- that is the gap I will return to in \S3b, and I now see this is because I equated ``moving the project forward'' with ``writing the next file,'' not with ``making sure two other people knew enough to write it without me.''

\textit{The cultural frame I brought, and what it cost me.} I grew up inside a Ukrainian engineering culture where the expected response to a blocker is to solve it yourself overnight and show the fix in the morning. That disposition served me on the IMX296 debug, the simulator build, and the complaint letter. It failed me on delegation: every time I chose ``I will just fix it myself'' I was, in hindsight, choosing ``nobody else on the team will learn this.'' The uncomfortable realisation \cite{hattonsmith1995} is that my greatest individual strength -- taking personal responsibility -- was simultaneously the team's largest single point of failure. That is the through-line of the whole report.

\textbf{1b. What I am most proud of.} \rtag{M16.d-own} Two things, in order.

\textit{The simulation framework, and the thesis behind it.} After three cancelled flight days (wk16, wk18, wk20 -- approximately 12 March, 19 March, 30~April) the team had accumulated zero seconds of real flight data on a course designed around a fly--test--tune iteration loop. I wrote the line that became the team's frame in the complaint letter to the course organiser \cite{complaintletter}: \textit{``The simulation framework we created ourselves is the only reason this project produced any verifiable engineering output at all.''} I still believe that is an accurate, not dramatic, statement. Without \texttt{simple\_simulator.py} -- GPS drift simulation, camera shake, CV throttle, TFLite backend, keyboard flight, live god-view -- the team would have had nothing to develop against between February and April. The simulator replaced approximately 50 physical test flights we never got to run. That it was not on any work-breakdown structure, and that I built it anyway \textit{before} the first flight day, is the single decision I am proudest of on this project, and it is the one decision that survives the ``would this hold up in a professional engineering review?'' test I keep trying to apply to my own work, in the sense of Sch\"on's reflection-in-action \cite{schon1983}.

\textit{The one debugging session that taught me how to be an engineer.} The IMX296 global-shutter camera was outputting what the \texttt{picamera2} library claimed was RGB888. The colours were wrong and detection failed. I spent hours on auto-white-balance modes, manual gains, gray-world correction, ISP tuning, even a kernel update. None of it worked. In the end I wrote a short loop that tried all six channel permutations and rendered them side by side. The sensor was actually outputting BGR. One-line fix. Ten files changed. The lesson was bigger than the bug: when a system lies about its interface, stop reading documentation and start enumerating possibilities. I wrote that rule into \texttt{MEMORY.md} so a future session would not repeat my day. I now see this is because the library's labelling had been an implicit authority I had not questioned -- and the deeper lesson is about authority in general, not just about colour channels. It is the episode I point to when I claim the project taught me something generalisable, rather than just a set of tools.

\textit{WWW / EBI.} What went well: autonomous initiative under external disruption, preference for runnable evidence over plans, willingness to abandon documentation when the artefact contradicted it. Even better if: I had paired on any of these wins with a teammate \textit{as I was doing them}, rather than handing them over as finished work afterwards. The proudness and the regret live in the same paragraph on purpose.

\textbf{1c. What I would focus differently.} \rtag{M16.d-own} Less solo building, earlier. I made the simulator excellent and the ground station browser-based specifically so that a teammate could pick up the mission logic in an afternoon, but I built it all by myself, which meant by the time anyone could ``pick it up,'' I was the only one who understood which bits were load-bearing. If I did this again I would pair-program from wk1, not wk20, even if the solo version would be faster. I would also refuse to take on the complaint letter single-handed: I wrote six iterations of it over four days [commits \texttt{26e3255 f3401e4 d946838 5af87d0 847518e 4b0fd91}] and that time should have been shared. I now see this is because I was, underneath, using ``taking it on alone'' as a proxy for ``caring enough'' -- and the two are not the same thing. The reflexive solo-fix disposition quietly subordinates honesty with colleagues to an ethic of personal delivery, which is exactly the substitution I now need to watch for.
